% ID: SEL-D0030-OB-E01
% ESTADO: validada. S031, cotejada con examen y criterios.
\textbf{a)} Al desarrollar, $(A+B)^2=A^2+AB+BA+B^2$.
La igualdad propuesta exigiría $AB=BA$. Sin embargo,
\[
AB=\begin{pmatrix}-2&-1\\2&-1\end{pmatrix},\qquad
BA=\begin{pmatrix}-1&2\\-1&-2\end{pmatrix}.
\]
Como difieren, la igualdad no se verifica.

\textbf{b)} Como $\det A=-2$, podemos multiplicar por $A^{-1}$
a la derecha. Así,
\[
X=(2B^t+I_2)A^{-1}
=\begin{pmatrix}5&0\\2&-1\end{pmatrix}
\begin{pmatrix}-1&0\\1/2&1/2\end{pmatrix}
=\boxed{\begin{pmatrix}-5&0\\-5/2&-1/2\end{pmatrix}}.
\]
Al sustituir en $XA$ se recuperan las filas $(5,0)$ y $(2,-1)$.
